% !TEX TS-program = pdflatex
% CV réseau LinkedIn — Secteur Bancaire — Infrastructure & Sécurité IT
\documentclass[a4paper,9pt]{article}
\usepackage[utf8]{inputenc}
\usepackage[T1]{fontenc}
\usepackage[scaled=0.90]{helvet}
\renewcommand{\familydefault}{\sfdefault}
\usepackage[left=0.8cm,right=0.8cm,top=0.3cm,bottom=0.25cm]{geometry}
\usepackage{enumitem}
\usepackage{titlesec}
\usepackage{xcolor}
\usepackage[hidelinks]{hyperref}
\definecolor{navy}{RGB}{23,54,93}
\pagestyle{empty}
\setlength{\parindent}{0pt}
\setlength{\parskip}{0pt}
\linespread{0.82}
\hypersetup{pdftitle={CV - Baha Eddine Belhaj Mouldi - Infrastructure et Securite IT Bancaire},pdfauthor={Baha Eddine Belhaj Mouldi},pdfsubject={Infrastructure, Securite, Systemes, Reseaux, Secteur Bancaire}}
\titleformat{\section}{\normalsize\bfseries\color{navy}}{}{0pt}{\MakeUppercase}[{\vspace{1pt}\color{navy}\hrule height 0.6pt}]
\titlespacing*{\section}{0pt}{1.5pt}{0.8pt}
\setlist[itemize]{leftmargin=1.0em,label=\textbullet,itemsep=0pt,topsep=0pt,parsep=0pt}
\newcommand{\cventry}[4]{\textbf{#1} \hfill \textbf{#4}\\[0.2pt]#2{\ifx\relax#3\relax\else, #3\fi}\par\vspace{0.3pt}}
\begin{document}
\begin{center}
{\LARGE\bfseries\color{navy} Baha Eddine Belhaj Mouldi}\\[2pt]
{\normalsize\color{navy} Ingénieur Infrastructure \& Sécurité IT \textbar\ Systèmes \textbar\ Réseaux \textbar\ Conformité}\\[3pt]
{\small Tunis, Tunisie \quad|\quad +216 55 128 982 \quad|\quad \href{mailto:bahaeddine.belhajmouldi@gmail.com}{bahaeddine.belhajmouldi@gmail.com}}\\[1pt]
{\small \href{https://www.linkedin.com/in/bahamouldi}{linkedin.com/in/bahamouldi} \quad|\quad \href{https://github.com/bahamouldi}{github.com/bahamouldi} \quad|\quad \href{https://bahamouldi.github.io/portfolio/}{bahamouldi.github.io/portfolio}}
\end{center}
\vspace{1pt}
\section{Profil}
Ingénieur en génie informatique (ENIT, mention Très Bien) avec une expérience concrète d'administration systèmes/réseaux et de sécurité applicative, acquise en travaux pratiques puis approfondie sur le terrain pendant un an de mission freelance et durant le stage de fin d'études chez Beehive Entreprises. Pratique de Windows Server / Active Directory (GPO, DHCP), Linux, virtualisation, sauvegarde (Veeam), haute disponibilité sur Kubernetes et conformité multi-référentiels (ISO 27001, NIST, PCI DSS, GDPR, SOC 2) au travers du projet BeeWAF. Intérêt marqué pour le secteur bancaire et financier, où la rigueur en matière de sécurité des systèmes d'information est centrale. Autonomie, rigueur et sens du résultat.
\section{Compétences clés}
\textbf{Systèmes \& Annuaire} : Windows Server, Active Directory (utilisateurs, groupes, GPO), Linux (Ubuntu, CentOS, Debian, Kali), durcissement système\\[0.3pt]
\textbf{Réseaux} : DHCP, DNS, VPN, Proxy, Firewall, segmentation réseau, TCP/IP\\[0.3pt]
\textbf{Virtualisation \& conteneurisation} : environnements virtualisés, Docker, Kubernetes (haute disponibilité, zéro downtime)\\[0.3pt]
\textbf{Sauvegarde \& continuité} : solutions de sauvegarde (Veeam), plans de continuité d'activité\\[0.3pt]
\textbf{Sécurité \& conformité} : OWASP Top 10, gestion des correctifs (CVE), PKI, MFA, IAM, ISO 27001, NIST, PCI DSS, GDPR, SOC 2\\[0.3pt]
\textbf{Supervision} : Prometheus, Grafana, ELK Stack, alerting temps réel\\[0.3pt]
\textbf{Scripting} : Python, Bash, PowerShell
\section{Expérience professionnelle}
\cventry{Ingénieur Sécurité \& Infrastructure --- Stagiaire PFE (mention Très Bien)}{Beehive Entreprises}{Tunisie}{02/2026 -- 06/2026}
\begin{itemize}
  \item Déploiement en production sur cluster Kubernetes (5 nœuds) : haute disponibilité, zéro downtime, supervision continue.
  \item Conformité continue sur 7 référentiels (OWASP, PCI DSS, GDPR, NIST, ISO 27001, CIS, SOC 2) ; gestion des correctifs (37 CVE).
\end{itemize}
\cventry{Spécialiste Infrastructure \& Sécurité, Freelance (1 an)}{Beehive Entreprises}{Tunisie}{01/2025 -- 12/2025}
\begin{itemize}
  \item Installation et administration Windows Server / Active Directory (GPO), Linux ; configuration DHCP, VPN, proxy, pare-feu.
  \item Virtualisation et sauvegarde (Veeam) ; audits de vulnérabilités sur 8 applications clientes.
\end{itemize}
\cventry{Stagiaire Cybersécurité --- Détection \& Réponse à incident SAP}{Streamlink}{Tunisie}{07/2025 -- 08/2025}
\begin{itemize}
  \item Supervision automatisée de transactions SAP critiques : détection $<$2 min, blocage automatisé $<$7 min.
\end{itemize}
\cventry{Chef de projet technique, Indépendant --- Équipe de 8 personnes}{VMate}{Tunisie}{01/2024 -- 06/2024}
\begin{itemize}
  \item Coordination d'équipe, planification, reporting.
\end{itemize}
\cventry{Stagiaire Sécurité Réseau}{Tunisie Telecom}{Tunisie}{07/2024}
\begin{itemize}
  \item Analyse de trafic réseau, durcissement, identification de 10+ vulnérabilités.
\end{itemize}
\section{Projets \& travaux pratiques}
\cventry{Infrastructure Windows Server / Active Directory \& Réseau Périmétrique}{Windows Server, AD, GPO, DHCP, VPN, Proxy, Firewall}{}{ENIT}
\begin{itemize}
  \item Administration AD, configuration DHCP/DNS, sécurisation périmétrique.
\end{itemize}
\cventry{BeeWAF --- Infrastructure de Sécurité en Production}{FastAPI, Docker, Kubernetes, Jenkins, ArgoCD, Prometheus, Grafana}{}{2026}
\begin{itemize}
  \item Cluster Kubernetes en production, zéro downtime, conformité continue sur 7 référentiels réglementaires.
\end{itemize}
\section{Formation}
\cventry{Diplôme d'Ingénieur en Génie Informatique --- Mention Très Bien}{École Nationale d'Ingénieurs de Tunis (ENIT)}{Tunisie}{09/2023 -- 06/2026}
\begin{itemize}
  \item Diplômé le 25/06/2026. Classé 65e sur 600+ au concours national d'entrée.
\end{itemize}
\cventry{Classes préparatoires Physique-Technologie}{Institut Préparatoire aux Études d'Ingénieurs, El Manar}{Tunisie}{09/2021 -- 06/2023}
\section{Certifications}
Réseaux (CCNA) --- Cisco (2025) ; Fondements du Deep Learning --- NVIDIA (2025) ; Machine Learning --- NVIDIA (2024) ; Git \& GitHub --- 365 Data Science (2024).
\section{Langues}
Arabe (Natif) \quad|\quad Français (Courant) \quad|\quad Anglais (B2)
\end{document}
